\newcommand{\Mod}[1]{\ (\mathrm{mod}\ #1)}
\DeclarePairedDelimiter{\ceil}{\lceil}{\rceil}
\newcommand\scalemath[2]{\scalebox{#1}{\mbox{\ensuremath{\displaystyle #2}}}}

\section{Shor's Factoring Algorithm}
% rename this section as you like, e.g. Applications of QFT, ...

In 1995, Peter Shor introduced a quantum algorithm that efficiently solves the integer factoring problem, i.e., decomposing large numbers into their prime factors in polynomial time. This was groundbreaking because factoring is believed to require exponential costs on a classical computer, which is why it became the basis of widely used cryptographic systems like RSA encryption\footnote{In fact, cryptographic standards are now evolving in response to the potential threat posed by quantum computers capable of running Shor’s algorithm.}. The implications of Shor's work are therefore enormous: Not only does it demonstrate a concrete case of practical quantum advantage, but it also provides a credible challenge to the strong Church-Turing thesis\footnote{in simple words, the thesis states that all efficient algorithms can be reduced to (probabilistic) Turing machines. Shor's Algorithm challenges the thesis since a quantum computer can efficiently factor primes, but to this day, there is no efficient classical algorithm. The prime-factoring problem is in fact a candidate for an $NP$-Intermediate problem -- conjectured to be neither $NP$-complete nor in $P$.}. This finding catalyzed the modern field of quantum computing, spurring both theoretical research and the development of actual quantum hardware. 

At its core, Shor’s algorithm is just a clever application of the quantum phase estimation \texttt{QPE} that includes the quantum Fourier transform \texttt{QFT} as a subalgorithm (see previous chapters). It relies on the observation that the factoring problem can be reduced to the order finding problem, which can be efficiently solved with \texttt{QPE}.

The quantum part of Shor's algorithm solves the order-finding problem. We will start with this and will then see how order-finding is used for integer factoring.

\subsection{Order-Finding}
Let us consider the following arithmetic problem.
\begin{def*}[Order-Finding]
Suppose an odd, composite number $N>1$ which is not a prime power, and an integer $x<N$, which is \textit{coprime} to $N$, i.e., $x$ and $N$ have no common factors. Determine the least positive integer $r$ such that ${x^r =1\Mod{N}}$. 
\end{def*}
In order words: At what order $r$ does the sequence ${x^1\Mod{N}},\ {x^2\Mod{N}}, ...$ reach the value $1$? As a simple example, consider the following: 
\begin{ex*}[$N=5$, $x=3$]
The coprime integers $N=5$ and $x=3$ give the sequence
\begin{align*}
 &3^1\Mod{5}= 3\Mod{5}= 3,\quad 3^2\Mod{5}= 9\Mod{5}= 4,\\
 &3^3\Mod{5}= 27\Mod{5}= 2, \quad 3^4\Mod{5}= 81\Mod{5}= 1, \quad \dots
\end{align*}
which tells us that the order of $3$ modulo $5$ is $4$. 
\end{ex*}

While this might seem easy to solve at first, note that for large integers, the value of $r$ can be very large, making an enumeration strategy like the one above highly impractical. In fact, order-finding is believed to be hard on classical computers, meaning that no known algorithm can solve it using only polynomial resources in $L\equiv\ceil{\log{N}}$, which is the number of bits encoding $N$. 

Instead, Shor proposed to solve this problem efficiently using quantum phase estimation (\texttt{QPE}) applied to the unitary operator
\begin{equation}
    U\ket{{y}} = \ket{{x}{y}\Mod{N}},
\end{equation}
where $\ket{{y}}$ denotes the basis state corresponding to the bit string encoding the integer $y$, with $0\leq y\leq N-1$. The states given by
\begin{equation}
    \left|u_s\right\rangle \equiv \frac{1}{\sqrt{r}} \sum_{k=0}^{r-1} \mathrm{e}^{-2 \pi \mathrm{i} k\frac{s}{r}}\ |x^k \Mod{N}\rangle,
\end{equation}
for integers $0\leq s\leq r-1$, are the eigenstates of $U$ with eigenvalues $\mathrm{e}^{2 \pi \mathrm{i} \frac{s}{r}}$. To show this, consider the following calculation:
\begin{equation}
\begin{aligned}
 U\left|u_s\right\rangle 
 & =\frac{1}{\sqrt{r}} \sum_{k=0}^{r-1}\ \mathrm{e}^{-2 \pi \mathrm{i} k\frac{s}{r}}\ |x^{k+1} \Mod{N}\rangle \\
& =\frac{1}{\sqrt{r}} \sum_{k=0}^{r-1}\ \mathrm{e}^{2 \pi \mathrm{i} \frac{s}{r}} \mathrm{e}^{-2 \pi \mathrm{i} (k+1)\frac{s}{r}}\ |x^{k+1} \Mod{N}\rangle \\
& \stackrel{(*)}{=} \mathrm{e}^{2 \pi \mathrm{i} \frac{s}{r}}\: \frac{1}{\sqrt{r}} \sum_{\tilde{k}=1}^r \mathrm{e}^{-2 \pi \mathrm{i} \tilde{k}\frac{s}{r}}\ |x^{\tilde{k}} \Mod{N}\rangle\qquad&\scalemath{0.8}{(*)\ \tilde{k}=k+1}\\
& \stackrel{(**)}{=}\mathrm{e}^{2 \pi \mathrm{i} \frac{s}{r}}\: \frac{1}{\sqrt{r}} \sum_{\tilde{k}=0}^{r-1} \mathrm{e}^{-2 \pi \mathrm{i} \tilde{k}\frac{s}{r}}\ |x^{\tilde{k}} \Mod{N}\rangle\qquad&\scalemath{0.8}{(**)\ \mathrm{periodicity\ w.r.t\ }r}\\
& =\mathrm{e}^{2 \pi \mathrm{i} \frac{s}{r}}\left|u_s\right\rangle
\end{aligned}
\end{equation}

This means that given the unitary $U$, an efficient procedure to implement controlled-$U^{2^j}$ operations\footnote{Note that this can be realized by a procedure called \textit{modular exponentiation} requiring $\mathcal{O}(L^3)$ gates.} for any $j$, and the corresponding eigenstate $\ket{u_s}$, we would be able to apply \texttt{QPE} to get the phase $\frac{s}{r}$ of the corresponding eigenvalue. Problem solved? Unfortunately, there is an issue: Preparing the state $\ket{u_s}$ requires prior knowledge of the value of $r$, which is precisely what we want to determine. However, we can circumvent this problem using the clever observation that the uniform superposition of all eigenstates satisfies\footnote{See Exercise sheet for a proof.}
\begin{equation}
\frac{1}{\sqrt{r}} \sum_{s=0}^{r-1}\left|u_s\right\rangle=|\mathbf{1}\rangle,
\end{equation}
which means we do not need to explicitly prepare $\ket{u_s}$. Instead, we can simply initialize the second register in the \texttt{QPE} algorithm in the easily prepared state $\ket{\mathbf{1}}$\footnote{Remember: $\ket{\mathbf{1}}$ denotes the computational basis state corresponding to the bit string $0\dots01$ of length $L$.}. As a result, when the \texttt{QPE} algorithm is executed, the first register collapses to an estimate of $\frac{s}{r}$ where $s\in\{0,\dots,r-1\}$ is chosen approximately uniformly at random. Since we are ultimately interested in the denominator $r$ rather than the particular value of $s$, this randomness does not affect the outcome. Further notice that by repeated applications of the unitary operator $\ket{U}$ to the state $\ket{1}$, we get 
\begin{equation}
    U^j\ket{1} = \ket{x^j\Mod{N}},
\end{equation}
which exactly reproduces the elements of the periodic sequence we are trying to analyze.

The quantum algorithm for order-finding is finally summarized as follows.

\begin{algo*}[\texttt{QPE} for Order-Finding]\phantom{...}\\
The following quantum circuit can be used to obtain an estimate for the phase $\frac{s}{r}$ given the state $\ket{1}$ and operator $U$ as defined above.
\begin{center}
\begin{quantikz}
\lstick{$\ket{0}$} & \qwbundle{n} & \gate{H^{\otimes n}} & \ctrl{1} & 
\gate{\texttt{QFT}^\dagger} & \meter{} \rstick{$\approx\frac{s}{r}$} \\
\lstick{$\ket{1}$} & \qwbundle{L} & \qw & \gate{ U^j } & 
\qw & \qw
\end{quantikz}
\end{center}
\end{algo*}
Note that the number of qubits in the first register determines the precision of our phase estimate. Specifically, if we choose
\begin{equation}
    n = 2L+1+\ceil[\Bigg]{\log(2+\frac{1}{2\epsilon})},
\end{equation}
then, for each $s\in\{0,\dots,r-1\}$, the above algorithm yields an estimate $\tilde{\varphi}$ of $\frac{s}{r}$ accurate to at least $2L+1$ bits, with success probability at least $\frac{1-\epsilon}{r}$. Once we have such an estimate $\tilde{\varphi}$, we can recover the exact value of the desired order parameter $r$ using the \textit{continued fractions} algorithm\footnote{see Tutorial for details on this procedure.}. With this final step, we have completed the reduction of order-finding to quantum phase estimation.

\subsection{Factoring}
The order-finding problem is computationally equivalent to the problem of factoring.
\begin{def*}[Factoring]
%Given a positive composite number $N$, what are the prime numbers -- when multiplied together -- equal it?
Given a positive composite number $N$, what are the prime numbers $X_1,X_2,\dots$ such that $\prod_k X_k = N$.
\end{def*}
This means that any efficient algorithm for order-finding can be straightforwardly transformed into an efficient factoring algorithm. The reduction from factoring to order-finding is based on two key properties:

\begin{theorem*}
Given an $L$-bit composite integer $N$, and suppose $x$ with $1\leq x\leq N$ satisfies
\begin{equation}
x^2=1\Mod{N},
\label{eq:modexp}
\end{equation}
but $x\neq \pm 1 \Mod{N}$. Then at least one of $\texttt{gcd}(x-1,N)$ and $\texttt{gcd}(x+1,N)$ is a non-trivial factor of $N$, computable using $\mathcal{O}(L^3)$ operations.
\end{theorem*}
This means that if we manage to find a non-trivial solution $x$ to \eqref{eq:modexp}, then we can compute a non-trivial factor of $N$, namely the greatest common divider ($\texttt{gcd}$) of $x\pm1$ and $N$.

\begin{theorem*}
Given an odd composite integer $N$ and its $m$ distinct prime factors. Let $x$ be a random integer chosen uniformly from $1,\dots, N-1$ such that $x$ is coprime to $N$, i.e., $\texttt{gcd}(x,N)=1$. Let $r$ be the order of $x$ modulo $N$. Then the probability of $r$ being even and 
\begin{equation}
    x^{\frac{r}{2}}\neq \pm 1 \Mod{N}
    \label{eq:modexp2}
\end{equation}
is at least $1-1/2^m$.
\end{theorem*}
This second theorem implies that if we chose a random number $y<N$, coprime to $N$, it is quite likely to have an order $r$ which is even and satisfies \eqref{eq:modexp2}. Linking this statement to the first property, we get that $x\equiv y^{\frac{r}{2}}\Mod{N}$ is a solution to \eqref{eq:modexp}.

Combining these considerations results in Shor's algorithm for integer factoring.
\begin{algo*}[Shor's Factoring Algorithm]\phantom{...}\\
Given a composite integer $N$, output a non-trivial factor of $N$.

\begin{algorithmic}
\STATE \textbf{if} $N \Mod{2} = 0$ \textbf{return} $2$ \hfill // check if $N$ is even
\STATE \textbf{if} $N = a^b$ for some $a \geq 1$, $b \geq 2$ \textbf{return} $a$ \hfill // check if $N$ is perfect power
\STATE \textbf{repeat}
\STATE $\quad x \gets$ random integer in $[1, N - 1]$
\STATE $\quad g \gets \texttt{gcd}(x, N)$
\STATE $\quad$\textbf{if} $g > 1$ \textbf{return} $g$
\STATE $\quad$\textbf{else}
\STATE $\quad\quad r \gets \texttt{QPE}$ for Order-Finding of $x$, $N$
\STATE $\quad\quad$\textbf{if} $r\Mod{2} = 0$ \textbf{and} $x^{r/2}\neq \pm 1\Mod{N}$ \textbf{do} $y \gets x^{r/2} \Mod{N}$
\STATE $\quad\quad\quad$\textbf{if} $y = 1$ \textbf{or} $y = N - 1$ \textbf{continue}
\STATE $\quad\quad\quad$\textbf{else}
\STATE $\quad\quad\quad\quad f_1 \gets \texttt{gcd}(y - 1, N)$
\STATE $\quad\quad\quad\quad f_2 \gets \texttt{gcd}(y + 1, N)$
\STATE $\quad\quad\quad\quad$\textbf{if} $1 < f_1 < N$ \textbf{return} $f_1$
\STATE $\quad\quad\quad\quad$\textbf{if} $1 < f_2 < N$ \textbf{return} $f_2$
\STATE $\quad\quad$\textbf{else} \textbf{continue}
\end{algorithmic}
\end{algo*}
The first two lines of preprocessing ensure that $N$ is an odd composite integer with more than one prime factor. This can be done using standard classical routines in $\mathcal{O}(1)$ and $\mathcal{O}\left(L^3\right)$ time. In the main step of the algorithm, we either immediately find a factor from a randomly selected element $x \in\{1, \ldots, N-1\}$ or we continue with the quantum order-finding subroutine from the previous section to compute the order $r$ of $x$ $\Mod{N}$. Using the above properties, we then attempt to extract a non-trivial factor of $N$ from $r$. This succeeds with probability at least $\frac{1}{2}$ for each iteration, provided $r$ is even and $x^{r / 2} \not \equiv \pm 1 \Mod{N}$.

Lastly note that all parts of this algorithm can be executed efficiently on a classical computer except for the order-finding subroutine, which is where the quantum advantage lies.
